% =====================================================
% ROTEIRO PRÁTICO 01 — PROGRAMAÇÃO PARA A INTERNET 2
% =====================================================

\section{Roteiro Prático 01 — Retomada e Transição: Bootstrap 5, jQuery e LocalStorage}

% =====================================================
% OBJETIVOS
% =====================================================

\begin{boxObjetivos}
\textbf{Objetivos da Aula}

Ao final desta aula o estudante será capaz de:

\begin{itemize}
\item Conectar a base de HTML5 semântico e CSS3 de \textbf{ProgWeb 1} com a evolução prática de \textbf{ProgWeb 2};
\item Compreender o alinhamento do semestre com o Projeto Pedagógico do Curso (\textbf{PPC} CONSUP 30/2020 — 120h);
\item Aplicar o \textbf{Bootstrap 5} para estilizar elementos de forma ágil (Exemplos 1 a 4: Grid, Formulários, Cards e Tabelas);
\item Empregar o \textbf{jQuery} explorando seus recursos funcionais, estéticos e dinâmicos de manipulação do DOM (Exemplos 5 a 7);
\item Persistir e recuperar dados de estado no navegador utilizando o \textbf{LocalStorage} (Exemplo 8).
\end{itemize}
\end{boxObjetivos}

% =====================================================
% INTRODUÇÃO
% =====================================================

\subsection{Introdução}

Bem-vindos à disciplina de \textbf{Programação para a Internet 2 (ProgWeb 2)}!

Neste primeiro encontro, construímos uma transição didática entre a base de **ProgWeb 1** e os novos recursos do semestre. Iniciamos de forma gradual: primeiro observando a transformação visual imediata dos elementos da página com o **Bootstrap 5** (Exemplos 1 a 4), em seguida manipulando o DOM e eventos com o **jQuery** divididos em recursos funcionais, estéticos e dinâmicos (Exemplos 5 a 7), finalizando com a persistência em **LocalStorage** (Exemplo 8). 

*Nota:* O estudo de consumo de APIs externas com a Fetch API e requisições HTTP assíncronas será aprofundado na Aula 02!

O trabalho é guiado pela \textbf{Apresentação Interativa de Transição} projetada em sala (\texttt{retomada-de-atividades-pi2/index.html}) e pela execução do projeto prático na pasta \texttt{exercicio-retomada/}.

% =====================================================
% PARTE 1
% =====================================================

\subsection{Parte 1 — Apresentação Interativa e Alinhamento com o PPC (25 minutos)}

\begin{boxAtividade}[{Atividade 1 -- Projeção dos Slides e Navegação pelos Exemplos 1 a 8}]

Acompanhamento da apresentação em slides dinâmicos (\texttt{retomada-de-atividades-pi2/index.html}):

\begin{enumerate}
\item \textbf{A Base de ProgWeb 1 \& PPC (Slides 2 e 3):} Recapitulação da estrutura HTML5 semântica e alinhamento com a ementa de 120h do curso;
\item \textbf{Bootstrap 5 em Prática (Slides 4 e 5):} Estilização direta de elementos com o Grid Responsivo (\textbf{Exemplo 1}), Formulários (\textbf{Exemplo 2}), Cards/Alertas (\textbf{Exemplo 3}) e Tabelas (\textbf{Exemplo 4});
\item \textbf{jQuery Funcional, Estético e Dinâmico (Slides 6, 7 e 8):} Seletores e leitura (\textbf{Exemplo 5}), Efeitos visuais \texttt{.fadeIn()} (\textbf{Exemplo 6}), Inserção dinâmica no DOM (\textbf{Exemplo 7}) e LocalStorage (\textbf{Exemplo 8});
\item \textbf{Guia Interativo de Código (Slide 9):} Consulta rápida das abas de código dos Exemplos 1 a 8 durante a aula.
\end{enumerate}

\end{boxAtividade}

% =====================================================
% PARTE 2
% =====================================================

\subsection{Parte 2 — Tarefa Prática: Dashboard com Bootstrap 5 e jQuery (50 minutos)}

\begin{boxAtividade}[{Atividade 2 -- Execução na Pasta exercicio-retomada/}]

Na pasta \texttt{exercicio-retomada/}, os estudantes deverão construir/completar a aplicação web:

\begin{enumerate}
\item \textbf{Estilização com Bootstrap 5:} Aplicar as classes de formulário e tabela do Bootstrap no \texttt{index.html} e observar o resultado visual;
\item \textbf{Seleção e Eventos com jQuery:} No arquivo \texttt{app.js}, capturar os inputs e os eventos de clique do formulário utilizando \texttt{\textdollar('\#id')};
\item \textbf{Inserção Dinâmica no DOM:} Adicionar linhas na tabela dinamicamente com \texttt{\textdollar('\#tabela-corpo').append()};
\item \textbf{Persistência em LocalStorage:} Salvar os dados do cliente (\texttt{JSON.stringify}) e permitir a exclusão de registros.
\end{enumerate}

\end{boxAtividade}

% =====================================================
% PARTE 3
% =====================================================

\subsection{Parte 3 — Desafio Extra: Filtro Dinâmico em Tempo Real (25 minutos)}

\begin{boxAtividade}[{Atividade 3 -- Filtro por Nome com jQuery}]

Como atividade de aprimoramento em sala:

\begin{enumerate}
\item Adicionar um campo de busca por nome acima da tabela (\texttt{\textdollar('\#filtro-nome')});
\item Implementar o evento \texttt{input} do jQuery para filtrar dinamicamente as linhas exibidas à medida que o usuário digita.
\end{enumerate}

\end{boxAtividade}

% =====================================================
% CHECKLIST
% =====================================================

\subsection{Checklist Final da Aula}

\begin{boxAtividade}[{Verificação Técnica}]

Antes de finalizar o encontro, verifique se seu projeto atende aos critérios de aceitação:

\begin{itemize}
\item A página utiliza as classes de formulário, botões e tabela do Bootstrap 5 (Exemplos 1 a 4);
\item O jQuery é utilizado para selecionar elementos do DOM (\texttt{\textdollar}) e reagir a eventos (Exemplos 5 a 7);
\item Os dados persistem no \texttt{localStorage} mesmo após recarregar a página (\texttt{F5}) (Exemplo 8);
\item O filtro por nome em tempo real e o botão de exclusão funcionam perfeitamente.
\end{itemize}

\end{boxAtividade}

% =====================================================
% REFLEXÃO TÉCNICA
% =====================================================

\subsection*{Reflexão Técnica}

A transição de ProgWeb 1 para ProgWeb 2 ganha ritmo ao combinarmos a simplicidade de estilização do Bootstrap 5 com a agilidade de manipulação do jQuery dividida em aspectos funcionais, estéticos e dinâmicos. Esse aprendizado gradual constrói a base sólida recomendada pelo PPC para avançarmos rumo à Fetch API e APIs REST na Aula 02!
